% ============================================================
% SECCIÓN 5: CRÍTICA Y EVALUACIÓN DE IMPACTO (ALINEACIÓN ABET)
% ============================================================

\section{Crítica y Evaluación de Impacto (Alineación ABET)}
\label{sec:impacto}

El \textit{ABET Student Outcome 4} exige que el ingeniero de computación reconozca responsabilidades profesionales y emita juicios informados basados en principios legales y éticos. El caso Tesla Autopilot revela tres niveles de responsabilidad que los tres marcos filosóficos iluminan desde ángulos distintos.

\textbf{Responsabilidad de diseño (nivel kantiano).} El ingeniero que codifica una función de pérdida para un sistema de VA no está tomando una decisión técnica neutral: está legislando una máxima de acción para millones de interacciones reales. El Código de Ética del IEEE exige que los ingenieros prioricen la seguridad pública sobre cualquier otro objetivo \citep{mishra2025moral}. Desde Kant, esta responsabilidad es categórica: no puede ser cedida a una lógica de mercado ni delegada a la ``objetividad'' del algoritmo. La parálisis deontológica no es una excusa para la inacción; es una señal de que la arquitectura de decisión debe rediseñarse hasta que ninguna de las opciones disponibles viole la dignidad individual.

\textbf{Responsabilidad de transparencia (nivel benthamita).} Bentham exige que el cálculo sea público. \citet{saetra2019tyranny} demuestran que los sistemas de Big Data que operan bajo funciones de optimización opacas generan una ``tiranía de la opinión percibida'': los individuos no pueden evaluar ni contestar los criterios bajo los cuales son clasificados, penalizados u optimizados. Para el VA, esto implica la obligación profesional de publicar los parámetros de la función de utilidad: ¿qué escenarios están cubiertos?, ¿bajo qué condiciones el sistema falla deliberadamente?, ¿quién asume el riesgo residual? La opacidad es, desde el benthamismo, una traición al principio democrático de que cada individuo debe poder calcular su propio riesgo.

\textbf{Responsabilidad de restricción (nivel milleano).} El principio del daño de Mill delimita el espacio de diseño éticamente admisible: el VA puede optimizar para el pasajero, pero no puede externalizar riesgo hacia terceros no consentientes sin su conocimiento. El principio del daño como restricción de diseño tiene una traducción técnica directa: los sistemas de seguridad que protegen a peatones y conductores circundantes no deben ser sujetos de optimización de rendimiento promedio. Son restricciones deontológicas dentro de un marco utilitarista de regla, exactamente como Mill propone \citep{mougan2024kantian}.

El análisis del caso Tesla revela una brecha de accountability que los tres marcos identifican desde ángulos distintos pero convergentes: cuando un sistema de IA mata, la responsabilidad se difumina entre el diseñador del algoritmo, el ingeniero de sistemas, la empresa, el regulador y el usuario. Esta difusión de responsabilidad es filosóficamente inaceptable bajo los tres marcos: Kant exige un agente racional que asuma la máxima del acto; Bentham exige un calculador que asuma la función de utilidad; Mill exige un legislador de reglas que asuma las consecuencias de las normas que impone. En ausencia de accountability clara, los tres sistemas éticos convergen en la misma conclusión: el diseño es moralmente irresponsable.
